\linksection{request}{The request Fixture}
Fixture functions can receive a \mono{request} object with attributes:

\begin{minted}[highlightlines=9-10,autogobble,escapeinside=||]{python}
    @pytest.fixture
    def myfix(request: pytest.FixtureRequest):
        request.function  # test function/method
        request.cls       # class of test
        request.instance  # class instance
        request.module    # module of test

        request.path      # path object of module
        request.node      # collection node
        request.config    # pytest config object

    def test_myfix(myfix):
        |\codeskip|
  \end{minted}

\vspace{-.5em}
\subsection{Getting Fixtures Dynamically}

With \py[escapeinside=||]{request.getfixturevalue(|\codeskip*|)}, we can get a fixture dynamically, to parametrize a test with fixtures:

\begin{minted}[autogobble]{python}
    @pytest.fixture
    def default_config():
        return Config("default.json")
\end{minted}
\begin{minted}[autogobble]{python}
    @pytest.fixture
    def debug_config():
        return Config("debug.json")
\end{minted}

\begin{minted}[escapeinside=||,autogobble]{python}
    @pytest.mark.parametrize("confname", [
        "default_config", "debug_config"])
    def test_configs(request, confname):
        config = |\highlightbox{request.getfixturevalue}|(confname)
        |\codeskip|
\end{minted}

Alternatives and (much!) more discussion:

\begin{itemize}
    \item pytest issue 349: \\ Using fixtures in \mono{pytest.mark.parametrize}
    \item The \mono{pytest-cases} plugin for more sophisticated parametrization
\end{itemize}
%\qrcode[height=4em]{https://github.com/pytest-dev/pytest/issues/349}

%\begin{tikzpicture}
%    \node [fixture, minimum width=2.5cm] (request) {\mono{request}};
%    \node [fixture, below=0.5 of request, minimum width=2.5cm] (myfix) {\mono{myfix}};
%    \node [test, below=0.75 of myfix, minimum width=2.5cm] (test) {\mono{test_myfix}};
%
%    \draw [->] (test) -- (myfix);
%    \draw [->] (myfix) -- (request);
%    \draw [-{to}, gray, rounded corners, label=right:foo]
%    (request.east) -- ([xshift=5mm]request.east) -- ([xshift=5mm]test.east) -- (test.east);
%\end{tikzpicture}

\subsection{Accessing Markers}

Via \py[escapeinside=||]{request.node.get_closest_marker(|\codeskip*|)}, you can get markers from test nodes.
As with e.g.\ \py[escapeinside=||]{@pytest.mark.parametrize(|\codeskip*|)},
you can pass
arguments to markers, and access them in your fixture:

\begin{minted}[escapeinside=||,autogobble]{python}
@pytest.fixture
def config(request: pytest.FixtureRequest):
    marker = request.node.get_closest_marker(
        "config_args")
    if marker is None:
        return Config()
    return Config(*marker.args, **marker.kwargs)
\end{minted}

The returned \mono{Mark} object has \mono{args} and \mono{kwargs} attributes with
the arguments passed to the marker. With

\begin{minted}[autogobble]{python}
@pytest.mark.config_args(
    "production.json", strict=True)
\end{minted}

you'll get: \hspace{6em}\raisebox{-.2em}{\textcolor{gray}{(single-element tuple)}}
\vspace{-.2em}

\begin{tabular}{lll}
    \py{marker.args}   & \textrightarrow{} & \py{("production.json",)} \\
    \py{marker.kwargs} & \textrightarrow{} & \py{{"strict": True}}
\end{tabular}

\subsection{Custom CLI Flags}
Using a plugin hook, you can also customize fixture behavior based on custom CLI flags:

\filenameheader{conftest.py}
\vspace{-.5em}

\begin{minted}[autogobble]{python}
    def pytest_addoption(parser: pytest.Parser):
        parser.addoption("--server-ip", type=str)

    @pytest.fixture
    def server_ip(request: pytest.FixtureRequest):
        return request.config.option.server_ip
\end{minted}
% \vspace{1em}

% \ex{}

% \begin{itemize}
%     \item Use this \mono{conftest.py}: In \mono{fixtures/cli_opt/}, create a \mono{test_server_ip.py} with a test using the \mono{server_ip} fixture
%     \item Invoke with different \mono{--server-ip} option values
% \end{itemize}

\vspace{.5em}
\filenameheadert{\$ pytest --help}
\vspace{-.5em}

\begin{minted}[escapeinside=||,autogobble]{text}
    |\codeskip|
    Custom options:
        |\highlightbox{--server-ip=SERVER\_IP}|
\end{minted}
